\begin{trapbox}
    \begin{description}[style=nextline, leftmargin=2.8em, labelsep=0.5em, itemsep=0.35em]
        \item[\textbf{\textcolor{CTrap}{易错点一（把候选根当成真正的根）}}]
            从 $p\mid a_0$、$q\mid a_n$ 就断言 $p/q$ 一定是根。
        \item[\textbf{\textcolor{CTrap}{正确理解（只有必要条件）}}]
            候选值必须代入验证。例如 $x^2+1$ 的可能有理根只有 $\pm1$，但 $f(1)=f(-1)=2\ne0$，实际上没有有理根。
        \item[\textbf{\textcolor{CTrap}{错因分析}}]
            有理根定理的方向是“\textbf{是根} $\Longrightarrow$ \textbf{满足整除}”，反方向不能随意使用。
    \end{description}

    \medskip
    \begin{description}[style=nextline, leftmargin=2.8em, labelsep=0.5em, itemsep=0.35em]
        \item[\textbf{\textcolor{CTrap}{易错点二（分子分母颠倒或未约分）}}]
            把分子拿去整除最高次项系数，或把 $2/4$ 直接当作 $p=2,q=4$。
        \item[\textbf{\textcolor{CTrap}{正确理解（首尾对应）}}]
            \textbf{分子看常数项，分母看最高次项系数}；先把根约成 $p,q$ 互质、$q>0$ 的形式。否则整除判断没有可靠依据。
        \item[\textbf{\textcolor{CTrap}{错因分析}}]
            证明里要用到 $p,q$ 互质；若分数没约分，就不能把 $q\mid a_np^n$ 进一步化为 $q\mid a_n$。
    \end{description}

    \medskip
    \begin{description}[style=nextline, leftmargin=2.8em, labelsep=0.5em, itemsep=0.35em]
        \item[\textbf{\textcolor{CTrap}{易错点三（常数项为零时套用 $p\mid a_0$）}}]
            多项式含有因子 $x$，却仍列“$0$ 的所有约数”去试有理根。
        \item[\textbf{\textcolor{CTrap}{正确理解（先处理零根）}}]
            $a_0=0$ 时先提取 $x$，确认 $x=0$，再对剩余多项式应用有理根定理；整除式中的 $p$ 不能取零。
        \item[\textbf{\textcolor{CTrap}{错因分析}}]
            整数 $0$ 不能作为除数；原定理采用 $a_0\ne0$ 的表述，正是为了先排除这种特殊情况。
    \end{description}

    \medskip
    \begin{description}[style=nextline, leftmargin=2.8em, labelsep=0.5em, itemsep=0.35em]
        \item[\textbf{\textcolor{CTrap}{易错点四（遗漏分母或重根）}}]
            只试 $\pm1,\pm2$ 等整数；找到一个根就停止，漏掉有理分数根或重复因式。
        \item[\textbf{\textcolor{CTrap}{正确理解（非首一必须试分数，降次还要继续）}}]
            如 $(2x-1)(x+1)$ 有根 $1/2$；如 $(x+1)^2$ 对应 $-1$ 是二重根。总次数要和原多项式核对。
        \item[\textbf{\textcolor{CTrap}{错因分析}}]
            试根发现的只是一个因式，不能代替对剩余商式的分析。
    \end{description}

    \medskip
    \begin{description}[style=nextline, leftmargin=2.8em, labelsep=0.5em, itemsep=0.35em]
        \item[\textbf{\textcolor{CTrap}{易错点五（把无有理根误说成不能分解）}}]
            认为多项式没有有理根，就一定不能写成两个整数系数多项式的乘积。
        \item[\textbf{\textcolor{CTrap}{正确理解（只能排除一次因式）}}]
            例如
            \[
              x^4+5x^2+4=(x^2+1)(x^2+4),
            \]
            没有有理根，但能分解成两个二次因式。
        \item[\textbf{\textcolor{CTrap}{错因分析}}]
            有理根与一次因式有关，不等于一切形式的多项式分解。
    \end{description}

    \medskip
    \begin{description}[style=nextline, leftmargin=2.8em, labelsep=0.5em, itemsep=0.35em]
        \item[\textbf{\textcolor{CTrap}{易错点六（忽略整数系数条件）}}]
            对带分数系数的多项式直接把分子分母套进“常数项、首项系数的约数”。
        \item[\textbf{\textcolor{CTrap}{正确理解（先化为整数系数）}}]
            如 $\frac12x^2-\frac34x+\frac14=0$，先乘 $4$，得到 $2x^2-3x+1=0$，再列候选有理根。
        \item[\textbf{\textcolor{CTrap}{错因分析}}]
            这一定理以整数系数为前提；分母不消去，无法进行整数的约数判断。
    \end{description}
\end{trapbox}
